\section{Designs and flight points}
\label{sec:flight}

\begin{table}[H]
\centering
\caption{The two 720\,mm noses. $L$: Haack length to the sharp apex; the \SI{4.5}{mm} spherical tip
removes $\delta$, so the curved part is $L - \delta$ long. Apogee: Python flight model calibrated on
Fluent (Section~\ref{sec:forces}).}
\label{tab:designs}
\small
\begin{tabular}{@{}lcc@{}}
\toprule
 & A (baseline) & B (backup)\\
\midrule
Profile & LV-Haack ($C = 1/3$) & Von K\'arm\'an ($C = 0$)\\
3U payload & inside the nose & not in the nose\\
$L$ / $\delta$ [mm] & \val{A-L} / \val{A-delta} & \val{B-L} / \val{B-delta}\\
Curved part $x_S$ + coupler $c$ [mm] & \val{A-xS} + \val{A-c} & \val{B-xS} + \val{B-c}\\
Fineness $L/D$ & \val{A-fineness} & \val{B-fineness}\\
Apogee, nominal / hot day [m] & \val{A-apogee} / \val{A-apogee-hot} & \val{B-apogee} / \val{B-apogee-hot}\\
\bottomrule
\end{tabular}
\end{table}

\paragraph{Flight.} The trajectory is the one of the Python design tool (report~2), with the nose of
each design. Two cases are given everywhere:
\begin{itemize}
\item \textbf{nominal}: ISA day (\SI{\val{pad-T}}{K} on the pad), propellant mass from the total impulse;
\item \textbf{hot day}: \SI{\val{pad-Thot}}{K} on the pad and the LVO masses. The rocket is lighter
and faster, so $q$ and the heating are higher.
\end{itemize}
Along the trajectory,
\begin{equation}
q = \tfrac12\,\rho_\infty V^2 = \tfrac12\,\gamma\,p_\infty M^2, \qquad
T_0 = T_\infty\left(1 + \tfrac{\gamma - 1}{2}\,M^2\right).
\end{equation}
The highest Mach number, dynamic pressure and stagnation temperature all occur at burnout
(Figure~\ref{fig:flight}).

\begin{figure}[H]
\centering
\includegraphics[width=\linewidth]{fig_flight}
\caption{(a) Mach number and (b) dynamic pressure along the flight, both designs, nominal and hot day.
Circles: the free stream of the Fluent runs.}
\label{fig:flight}
\end{figure}

\paragraph{Fluent points.} The three Fluent runs use the free stream of design~A's nominal trajectory
at these instants (Table~\ref{tab:points}); design~B is run at the same points, so that A and B differ
by shape only (B's own burnout is at $M = \val{B-own-burnout-M}$). Values at other times are interpolated
in $M$ (Section~\ref{sec:pressure}).

\begin{table}[H]
\centering
\caption{Free stream of the Fluent runs. Reynolds number per metre: $\rho_\infty V/\mu_\infty$.}
\label{tab:points}
\small
\begin{tabular}{@{}lccc@{}}
\toprule
 & Transonic & Supersonic & Burnout\\
\midrule
Time $t$ [s] & \val{pt-M095-t} & \val{pt-M120-t} & \val{pt-M147-t}\\
Altitude $z$ [m] & \val{pt-M095-z} & \val{pt-M120-z} & \val{pt-M147-z}\\
Mach number $M$ & \val{pt-M095-M} & \val{pt-M120-M} & \val{pt-M147-M}\\
Static pressure $p_\infty$ [Pa] & \val{pt-M095-p} & \val{pt-M120-p} & \val{pt-M147-p}\\
Static temperature $T_\infty$ [K] & \val{pt-M095-T} & \val{pt-M120-T} & \val{pt-M147-T}\\
Dynamic pressure $q$ [kPa] & \val{pt-M095-q} & \val{pt-M120-q} & \val{pt-M147-q}\\
Reynolds number per metre [$10^6$/m] & \val{pt-M095-Re} & \val{pt-M120-Re} & \val{pt-M147-Re}\\
\bottomrule
\end{tabular}
\end{table}

\paragraph{Load cases.} Three instants bound the aerodynamic loads of the sealed nose:
\begin{itemize}
\item \textbf{LC1, burnout}: highest $M$, $q$, inward pressure at the tip, axial force, normal force per
degree and heating;
\item \textbf{LC2, transonic}, $M = 0.95$ on the ascent: the shock system changes, lower $q$;
\item \textbf{LC3, apogee}: $q \approx 0$ and the lowest outside pressure, so the largest outward
(burst) pressure on the sealed nose.
\end{itemize}
Recovery (Req\_07) and the inertial loads at the joint are not aerodynamic and are left to ST.
